\documentclass[aspectratio=169,12pt]{beamer}

% ============================================================
%  Apresentação: Detecção de Phishing em URLs com
%  Classificadores Clássicos — Um Estudo Comparativo
%  Divisão: 3 apresentadores (blocos marcados com \parte{})
% ============================================================

\usepackage[utf8]{inputenc}
\usepackage[brazil]{babel}
\usepackage{booktabs}
\usepackage{tikz}
\usetikzlibrary{shapes.geometric,arrows.meta,positioning}

% ---- Tema visual limpo ----
\usetheme{Madrid}
\usecolortheme{whale}
\useinnertheme{circles}
\setbeamertemplate{navigation symbols}{}
\setbeamertemplate{caption}{\raggedright\insertcaption\par}
\setbeamerfont{frametitle}{size=\large}
\setbeamertemplate{frametitle}[default][center]

% Rodapé enxuto e neutro (sem apresentador)
\setbeamertemplate{footline}{%
  \hbox{%
    \begin{beamercolorbox}[wd=\paperwidth,ht=2.5ex,dp=1.2ex,leftskip=1em,rightskip=1em]{author in head/foot}%
      \usebeamerfont{author in head/foot}\insertshorttitle\hfill
      \insertframenumber/\inserttotalframenumber
    \end{beamercolorbox}%
  }%
}

% Separador de seção neutro (sem apresentador); alimenta o sumário
\newcommand{\parte}[1]{%
  \section{#1}
  \begingroup
  \setbeamercolor{background canvas}{bg=structure.fg!12}
  \begin{frame}[plain]
    \vfill
    \begin{center}
      {\usebeamercolor[fg]{structure}\Large\textbf{#1}}
    \end{center}
    \vfill
  \end{frame}
  \endgroup
}

% Dados reais do experimento (espelham resultados/macros.tex)
\newcommand{\nTotal}{11.055}
\newcommand{\nFeatures}{30}
\newcommand{\nPhishing}{6.157}
\newcommand{\nLegitimo}{4.898}
\newcommand{\pctPhishing}{55,7\%}
\newcommand{\melhorFOne}{0,974}
\newcommand{\melhorAUC}{0,998}
\newcommand{\melhorAcuracia}{0,974}
\newcommand{\naiveBayesFOne}{0,561}
\newcommand{\friedmanChi}{56,27}
\newcommand{\topDoisPct}{56,7\%}

\title[Detecção de Phishing em URLs]{Detecção de Phishing em URLs com\\Classificadores Clássicos}
\subtitle{Um Estudo Comparativo}
\author[Dias, Braz e Nascimento]{Vinícius M. Dias \and João M. M. Braz \and Pedro G. A. do Nascimento}
\institute[UESPI]{Ciência da Computação — Universidade Estadual do Piauí (UESPI)}
\date{Reconhecimento de Padrões}

\begin{document}

% ===================== CAPA =====================
\begin{frame}[plain]
  \titlepage
\end{frame}

% ===================== SUMÁRIO =====================
\begin{frame}{Sumário}
  \tableofcontents
\end{frame}

% ============================================================
% ================  PARTE 1 — VINÍCIUS  ======================
% ============================================================
\parte{Introdução e Trabalhos Relacionados}

\begin{frame}{O problema: phishing}
  \begin{itemize}
    \item Phishing é engenharia social: URLs fraudulentas imitam serviços
          legítimos para roubar credenciais ou instalar malware.
    \item \textbf{Escala do problema:} só em 2023, a APWG registrou cerca de
          \textbf{cinco milhões} de ataques — o maior total anual já registrado.
    \item Triagem manual é inviável $\Rightarrow$ filtros automáticos com
          \textbf{aprendizado de máquina} tornaram-se essenciais em browsers,
          e-mail e proxies.
  \end{itemize}
  \vfill
  \begin{block}{Pergunta do trabalho}
    Qual classificador clássico é mais eficaz para detectar phishing em URLs,
    e a diferença entre eles é estatisticamente significativa?
  \end{block}
\end{frame}

\begin{frame}{O que este trabalho faz}
  \begin{itemize}
    \item Compara \textbf{sete classificadores clássicos}:
    \begin{itemize}
      \item Regressão Logística, Naive Bayes, KNN, Árvore de Decisão,
      \item Random Forest, SVM (kernel RBF), Gradient Boosting.
    \end{itemize}
    \item Base: \textbf{UCI Phishing Websites} — \nTotal{} amostras,
          \nFeatures{} atributos.
    \item Avaliação com validação cruzada de 10 folds \textbf{e} teste
          estatístico de Friedman.
  \end{itemize}
\end{frame}

\begin{frame}{Trabalhos relacionados}
  \begin{itemize}
    \item \textbf{Mohammad et al. (2012)} — propuseram o dataset UCI, definindo e
          extraindo automaticamente os \nFeatures{} atributos.
    \item \textbf{Sahingoz et al. (2019)} — compararam sete algoritmos; Random
          Forest venceu (97,98\% de acurácia), na mesma faixa do nosso resultado.
    \item \textbf{Jain e Gupta (2018)} — classificação por hiperlinks do lado do
          cliente (98,4\%), sem carregar a página.
    \item \textbf{Verma e Das (2017)} — atributos léxicos da string da URL já são
          altamente discriminativos.
  \end{itemize}
\end{frame}

% ============================================================
% ================  PARTE 2 — JOÃO MARCELLO  =================
% ============================================================
\parte{Fundamentação e Metodologia}

\begin{frame}{Os sete classificadores}
  \small
  \begin{tabular}{@{}ll@{}}
    \toprule
    \textbf{Modelo} & \textbf{Ideia central} \\
    \midrule
    Regressão Logística & Fronteira linear; probabilidade sigmoidal \\
    Naive Bayes Gaussiano & Assume independência e normalidade \\
    KNN ($k=5$) & Vota nos vizinhos mais próximos \\
    Árvore de Decisão & Partições por impureza de Gini \\
    Random Forest & Ensemble de árvores (reduz variância) \\
    SVM (kernel RBF) & Margem máxima em espaço transformado \\
    Gradient Boosting & Árvores aditivas que corrigem resíduos \\
    \bottomrule
  \end{tabular}
  \vfill
  \footnotesize Modelos sensíveis à escala (RegLog, NB, KNN, SVM) recebem
  padronização; os baseados em árvore, não.
\end{frame}

\begin{frame}{O dataset}
  \begin{itemize}
    \item \textbf{UCI Phishing Websites}: \nTotal{} amostras, \nFeatures{}
          atributos categóricos discretos.
    \item Maioria \textbf{binários} $\{-1,1\}$ (ausência/presença);
          alguns \textbf{ternários} $\{-1,0,1\}$ (intensidade).
    \item \nPhishing{} phishing (\pctPhishing) e \nLegitimo{} legítimas
          — leve desbalanceamento, sem valores faltantes.
    \item Atributos extraídos de URL, DNS, HTML e JavaScript.
  \end{itemize}
\end{frame}

\begin{frame}{Pipeline experimental}
  \centering
  \scriptsize
  \begin{tikzpicture}[
    node distance=3.2mm,
    etapa/.style={rectangle, rounded corners, draw, thick,
      text width=6.2cm, align=center, minimum height=6mm, inner sep=2pt},
    seta/.style={-{Stealth[length=1.6mm]}, thick}
  ]
    \node[etapa, fill=black!5] (d) {Dataset UCI (\nTotal{} amostras, \nFeatures{} atributos)};
    \node[etapa, fill=blue!8, below=of d] (s) {Divisão treino/teste (80/20, estratificada)};
    \node[etapa, fill=blue!8, below=of s] (p) {Padronização (modelos sensíveis à escala)};
    \node[etapa, fill=violet!12, below=of p] (cv) {Validação cruzada estratificada de 10 folds};
    \node[etapa, fill=violet!12, below=of cv] (t) {Treino dos sete classificadores};
    \node[etapa, fill=teal!14, below=of t] (a) {Avaliação no teste (acurácia, F1, AUC-ROC)};
    \node[etapa, fill=teal!14, below=of a] (st) {Testes estatísticos (Friedman + Nemenyi)};
    \draw[seta] (d)--(s); \draw[seta] (s)--(p); \draw[seta] (p)--(cv);
    \draw[seta] (cv)--(t); \draw[seta] (t)--(a); \draw[seta] (a)--(st);
  \end{tikzpicture}
\end{frame}

\begin{frame}{Como medimos e comparamos}
  \begin{itemize}
    \item \textbf{Protocolo:} 80/20 estratificado, \texttt{random\_state=42}
          (reprodutível); validação cruzada de 10 folds no treino.
    \item \textbf{Métricas:} acurácia, F1 macro e AUC-ROC.
    \item \textbf{Teste estatístico:} Friedman (as diferenças são reais?)
          + pós-teste de Nemenyi (quais modelos diferem?).
    \item \textbf{Ferramentas:} scikit-learn, SciPy, scikit-posthocs.
  \end{itemize}
  \vfill
  \begin{block}{Por que teste estatístico?}
    Diferenças pequenas de F1 podem ser ruído. O Friedman diz se há diferença
    \emph{significativa}, não apenas numérica.
  \end{block}
\end{frame}

% ============================================================
% ================  PARTE 3 — PEDRO GABRYEL  =================
% ============================================================
\parte{Resultados e Conclusão}

\begin{frame}{Resultados: desempenho dos modelos}
  \centering
  \scriptsize
  \begin{tabular}{lcccc}
    \toprule
    \textbf{Classificador} & \textbf{F1 (CV)} & \textbf{F1 (teste)} & \textbf{AUC-ROC} & \textbf{Tempo (s)} \\
    \midrule
    Regressão Logística & 0,926 & 0,927 & 0,981 & 0,01 \\
    Naive Bayes & 0,561 & 0,561 & 0,970 & 0,00 \\
    KNN ($k=5$) & 0,940 & 0,949 & 0,987 & 0,00 \\
    Árvore de Decisão & 0,961 & 0,971 & 0,980 & 0,01 \\
    \textbf{Random Forest} & \textbf{0,970} & \textbf{0,974} & \textbf{0,998} & 0,17 \\
    SVM RBF & 0,948 & 0,951 & 0,989 & 2,14 \\
    Gradient Boosting & 0,952 & 0,956 & 0,994 & 0,75 \\
    \bottomrule
  \end{tabular}
  \vfill
  \normalsize
  \textbf{Random Forest} vence: F1 = \melhorFOne, AUC-ROC = \melhorAUC,
  com baixa variância e treino em fração de segundo.
\end{frame}

\begin{frame}{Dois destaques da análise}
  \begin{columns}[T]
    \begin{column}{0.5\textwidth}
      \textbf{Naive Bayes despenca}
      \begin{itemize}
        \item F1 = \naiveBayesFOne{} (muito abaixo).
        \item AUC-ROC até competitivo (0,970), mas probabilidades mal
              calibradas $\Rightarrow$ limiar 0,5 falha.
        \item Atributos discretos violam a hipótese gaussiana.
      \end{itemize}
    \end{column}
    \begin{column}{0.5\textwidth}
      \textbf{Árvore de Decisão}
      \begin{itemize}
        \item 2º melhor F1 (0,971)\ldots
        \item \ldots mas o \emph{menor} AUC dos bons (0,980).
        \item Decide bem no limiar padrão, mas gradua mal as probabilidades.
      \end{itemize}
    \end{column}
  \end{columns}
\end{frame}

\begin{frame}{Teste estatístico e atributos}
  \begin{itemize}
    \item \textbf{Friedman:} $\chi^2 = \friedmanChi{}$ ($p < 0,001$) — as
          diferenças são significativas.
    \item \textbf{Nemenyi:} só o Naive Bayes difere dos demais; os outros seis
          são estatisticamente equivalentes — todos viáveis.
    \item \textbf{Atributos mais discriminativos:} \texttt{sslfinal\_state}
          (SSL) e \texttt{url\_of\_anchor} (âncoras externas) somam \topDoisPct{}
          da importância.
  \end{itemize}
  \vfill
  \footnotesize No artigo: Figura 2 (curvas ROC) e Figura 3 (importância dos
  atributos) detalham estes pontos.
\end{frame}

\begin{frame}{Conclusão}
  \begin{itemize}
    \item Random Forest foi o melhor: F1 = \melhorFOne{}, AUC-ROC = \melhorAUC{},
          rápido e estável.
    \item As diferenças entre os modelos são estatisticamente significativas
          (Friedman); o Naive Bayes é o destoante.
    \item SSL e proporção de âncoras externas concentram o poder discriminativo.
  \end{itemize}
  \vfill
  \begin{block}{Limitação e trabalhos futuros}
    Os atributos são extraídos manualmente (implica latência). Próximo passo:
    atributos léxicos brutos da URL, sem carregar a página.
  \end{block}
\end{frame}

\begin{frame}[plain]
  \vfill
  \begin{center}
    {\Large\textbf{Obrigado!}}\\[1em]
    {\large Perguntas?}\\[2em]
    \small Vinícius M. Dias \quad\textbullet\quad João M. M. Braz
    \quad\textbullet\quad Pedro G. A. do Nascimento\\
    \small UESPI — Reconhecimento de Padrões
  \end{center}
  \vfill
\end{frame}

\end{document}
